\section{The physical half-wave and frequency summation}
\label{sec:phys}

We now deduce Theorem~\ref{thm:main} from the wave assertion of
Theorem~\ref{thm:round}.  We use the negative half-wave
\[
 Wg(x,\tau)=(2\pi)^{-3}\int_{\R^3}
 e^{i(x\cdot\xi-\tau|\xi|)}\widehat g(\xi)\,d\xi.
\]
Complex conjugation, followed by the change of variable $\xi\mapsto-\xi$,
interchanges the two signs.  The same conjugation commutes with every
real radial Sobolev multiplier.

There are two points to verify before the round-cone estimate applies to
the physical wave: the packet state must have the required regularity on
the original time slice, and its terminal masses must control the physical
$L^3$ norm. We prove both at one frequency scale, then sum over spatial
pieces, amplitude levels, and frequency annuli.

\subsection{The estimate on one annulus}

For $0<H\le1$, let $a_H$ be a smooth multiplier supported where
$c\le H|\xi|\le C$.  All derivatives of $a_H(H^{-1}\cdot)$ are assumed
uniformly bounded.  Write
\begin{equation}
 W_Hg(x,\tau)=(2\pi)^{-3}\int
 e^{i(x\cdot\xi-\tau|\xi|)}a_H(\xi)\widehat g(\xi)\,d\xi.
 \label{eq:phys-annular-wave}
\end{equation}
The constants below may depend on the fixed annulus and on these
smoothness bounds.  They do not depend on $H$ or on $g$.

\begin{proposition}[Annular estimate]
\label{prop:phys-annular}
For every $\delta>0$,
\begin{equation}
 \norm{W_Hg}_{L^3(\R^3\times[1,2])}
 \le C_\delta H^{-\delta}\norm{g}_{L^3(\R^3)}.
 \label{eq:phys-annular-estimate}
\end{equation}
\end{proposition}

It suffices to prove the proposition for dyadic $H$: any $H$ is
comparable to a dyadic parameter, and replacing it by that parameter
changes the fixed support and derivative bounds by fixed constants.
We prove the proposition in several steps.  A finite smooth angular
partition and rotations reduce its proof to multipliers for which
\begin{equation}
 \rho=\frac{|\xi|+\xi_3}{2},\qquad
 A=\frac{\xi_\perp}{2\rho},\qquad
 c_0\le H\rho\le C_0,\quad |A|\le A_0.
 \label{eq:phys-patch}
\end{equation}
Such a partition exists because each direction has a neighborhood on
which, after a rotation, $|\xi|+\xi_3$ is comparable to $|\xi|$.

We can also restrict attention to bounded input and output regions.
Indeed, if $K_{H,\tau}$ is the convolution kernel of
Equation~\eqref{eq:phys-annular-wave}, then, for every $N$, there is a
fixed $R$ such that
\begin{equation}
 |K_{H,\tau}(x)|\le C_N H^N(1+|x|)^{-N}
 \quad (|x|\ge R,\ 1\le\tau\le2).
 \label{eq:phys-far-kernel}
\end{equation}
To see this, substitute $\xi=H^{-1}\omega$.  The gradient of the phase
in $\omega$ is $H^{-1}(x-\tau\omega/|\omega|)$, whose magnitude is
comparable to $H^{-1}(1+|x|)$ in the indicated region.  Repeated
integration by parts absorbs the initial factor $H^{-3}$ and proves
the estimate.  Partition output space into unit cubes $Q_j$, and let
$Q_j^*$ be fixed enlargements large enough that $|x-z|\ge R$ whenever
$x\in Q_j$ and $z\notin Q_j^*$.  On each $Q_j$ replace $g$ by
$g\mathbf1_{Q_j^*}$.  The error is pointwise bounded by the convolution
of $|g|$ with the right side of
Equation~\eqref{eq:phys-far-kernel}.  Its $L^3$ norm is $O_N(H^N)\norm g_3$
by Young's inequality.  The cubes $Q_j^*$ have bounded overlap.
Consequently a uniform local estimate can be cubed and summed over
$j$.  Translation reduces these local estimates to input supported
in one fixed compact set $K$ and output in one fixed bounded set.

\subsection{The null-coordinate initial state}

Set
\begin{equation}
 t=\tau+x_3,\qquad b=(x_1,x_2),\qquad y=\tau-x_3.
 \label{eq:phys-null-coordinates}
\end{equation}
The absolute spacetime Jacobian is
$dt\,db\,dy=2\,d\tau\,dx$.  In the frequency variables
$(\eta,\rho)=(\xi_\perp,(|\xi|+\xi_3)/2)$, one has
\begin{equation}
 \xi_3=\rho-\frac{|\eta|^2}{4\rho},\quad
 |\xi|=\rho+\frac{|\eta|^2}{4\rho},\quad
 J(\eta,\rho):=\frac{\partial\xi_3}{\partial\rho}
       =1+\frac{|\eta|^2}{4\rho^2}=1+|A|^2.
 \label{eq:phys-frequency-change}
\end{equation}
Thus the same function, denoted by $v$, is exactly
\begin{equation}
 v(t,b,y)=(2\pi)^{-3}\int
 e^{i(b\cdot\eta-y\rho-t|\eta|^2/(4\rho))}
 J(\eta,\rho)a_H(\xi(\eta,\rho))
 \widehat g(\xi(\eta,\rho))\,d\eta\,d\rho.
 \label{eq:phys-round-state}
\end{equation}
It is therefore a free round-cone wave.  Its energy is independent of
$t$ and satisfies
\begin{equation}
 e:=\norm{v(t)}_2^2
 =(2\pi)^{-3}\int J(\eta(\xi),\rho(\xi))
          |a_H(\xi)\widehat g(\xi)|^2\,d\xi
 \le C\norm g_2^2.
 \label{eq:phys-energy}
\end{equation}
In particular the map from $g$ to any plateau-analyzed state
$V_1v(T)$ has uniformly bounded $L^2$ norm.  This assertion remains
true after replacing $g$ by its restriction to any measurable subset
of $K$.

The bounded physical output region is covered by finitely many unit
intervals $[T,T+1]$ in $t$, with bounded $T$.  We work on one such
interval.  Write $B=\norm g_\infty$, and suppose first that
$g\in L^\infty$ is supported in $K$.  If $B=0$ or $e=0$, the asserted
estimate is immediate.  The frequency support in
$(\eta,-\rho)$ has volume $O(H^{-3})$, so Cauchy--Schwarz in Fourier
space gives
\begin{equation}
 \norm{v(t)}_\infty\le C H^{-3/2}\sqrt e.
 \label{eq:phys-bernstein}
\end{equation}
If $e<H^{10}B^2$, then
\begin{equation}
 \int_T^{T+1}\int |v|^3
 \le C H^{-3/2}e^{3/2}
 \le C B\norm g_2^2.
 \label{eq:phys-small-energy}
\end{equation}
We may therefore assume
\begin{equation}
 H^{10}B^2\le e\le C\norm g_2^2\le C_K B^2.
 \label{eq:phys-energy-range}
\end{equation}

\subsection{Input regularity on the physical initial slice}

The physical initial slice $\tau=0$ is the spacelike plane $y=-t$.
We verify the full regularity family on this plane, including the
nonround tests. After a kernel truncation uniform over all rows, let
$E\subset K$ comprise the physical initial points joined to a test by
retained kernel pairs. Equation~\eqref{eq:phys-energy} then bounds the
test's row mass by $CB^2|E|$, up to an error negligible relative to
$\norm g_2^2$. We therefore need a uniform kernel-localization estimate
and a bound for $|E|$ by the test weight $r^4f^{1-\eta}$, allowing an
arbitrarily small power loss.

\begin{lemma}[Uniform packet localization from the initial slice]
\label{lem:phys-initial-kernel}
Fix $\gamma_0>0$ and put $\Lambda=H^{-\gamma_0}$.  The operator
$g\mapsto V_1v(T)$ can be truncated to the following relation between
an input point $z=(z_\perp,z_3)\in K$ and a row $(A,b,y)$:
\begin{align}
 |b-z_\perp-(T-z_3)A|&\le C\Lambda\sqrt H,\notag\\
 |y+z_3-(T-z_3)|A|^2
       -2A\cdot(b-z_\perp-(T-z_3)A)|&\le C\Lambda H.
 \label{eq:phys-predecessor-relation}
\end{align}
For every $N$, the discarded operator has norm at most $C_NH^N$
from $L^2(K)$ to the row space.  The constant is independent of any
regularity test subsequently imposed on the output rows.
The kernel vanishes unless $|A|\le A_0+C_\psi\sqrt H$, where
$C_\psi$ is a radius containing the support of the fixed angular
window $\psi$.
\end{lemma}

\begin{proof}
Insert the Fourier representation of $\widehat g$ into
Equation~\eqref{eq:phys-round-state} and the plateau analysis.  The
kernel has phase
\[
 (b-z_\perp)\cdot\eta-(y+z_3)\rho
                  -(T-z_3)\frac{|\eta|^2}{4\rho}.
\]
Its amplitude is supported in $H\rho\asymp1$ and
$|\eta/(2\rho)-A|\le C_\psi\sqrt H$.
Together with $|\eta/(2\rho)|\le A_0$ from
Equation~\eqref{eq:phys-patch}, this proves the stated angular
support bound.  On making the substitutions
$\eta=2\rho A+H^{-1/2}\omega$ and $\sigma=H\rho$, the frequency
Jacobian is $H^{-2}$, and the amplitude has uniformly bounded
derivatives on a fixed compact set.  The quadratic remainder in the
phase also has bounded derivatives, because $T-z_3$ is bounded.
Fourier integration by parts therefore gives
\[
 |\mathcal K_T(A,b,y;z)|\le C_L H^{-2}
 \left(1+\frac{|\Delta b|}{\sqrt H}
              +\frac{|\Delta y|}{H}\right)^{-L},
\]
where $\Delta b$ and $\Delta y$ are the two differences in
Equation~\eqref{eq:phys-predecessor-relation}.  Only bounded $A$ occur.
The row measure at length $1$ is $H^{-1}\,dA\,db\,dy$.
For fixed $A,z$, integration in $(b,y)$ costs the packet volume
$H^2$.  Squaring the kernel and then integrating over all rows and
$z\in K$ thus gives a Hilbert--Schmidt bound of size at most a fixed
power of $H^{-1}$.  Outside the displayed relation, the decay factor
adds an arbitrarily large inverse power of $\Lambda$.
Since $\gamma_0$ is fixed, taking $L$ sufficiently large proves any
prescribed $H^N$ bound.  This estimate was made over all output rows,
which proves the asserted uniformity.
\end{proof}

\begin{lemma}[Volume of physical predecessors]
\label{lem:phys-predecessor-volume}
Fix $0<\eta<1$.  Consider any round-cone regularity test at time $T$
and length $1$, of angular width $r$ and relative thickness $f$,
where $0<f\le1$ and $rf\ge\sqrt H$; for a round test set $f=1$.
Let $E$ be the set of $z\in K$ related by
Equation~\eqref{eq:phys-predecessor-relation} to some row of the test
whose angular coordinate satisfies $|A|\le A_0+C_\psi\sqrt H$.
There is a fixed $C_1$ such that
\begin{equation}
 |E|\le C\Lambda^{C_1}r^4 f^{1-\eta}.
 \label{eq:phys-predecessor-volume}
\end{equation}
The estimate is uniform in the center, $r$, and $f$ of the test.
\end{lemma}

\begin{proof}
Use its normalized coordinates
\[
 u=t-T,\quad x=\frac{b-b_*-uA_*}{r},\quad
 w=\frac{y-y_*-2A_*\cdot(b-b_*)+u|A_*|^2}{r^2},\quad
 V=\frac{A-A_*}{r}.
\]
At $u=0$, the test has $|x|,|w|,|V|\le1$.  The prediction at $T$
from a compatible input point differs from its tested row by
\[
 |\delta x|\le C\Lambda\frac{\sqrt H}{r}\le C\Lambda f,
 \qquad
 |\delta w|\le C\Lambda\left(\frac H{r^2}
                      +|V|\frac{\sqrt H}{r}\right)
                  \le C\Lambda f.
\]
Backflow to $t=z_3$ has bounded duration.  Consequently the
compatible normalized coordinates are bounded by a fixed power of
$\Lambda$.  For a nonround test, write its polynomial as
\[
 F=c+du+\alpha\cdot x+\beta w+k(uw-|x|^2),\qquad
 G=F_u+V\cdot F_x+|V|^2F_w.
\]
The coefficients are bounded, their discriminant is $1$, $G$ is
constant on rays, and $F$ is affine on rays.  Expanding at the
tested row and then backflowing shows that, at a compatible input
point,
\begin{equation}
 |F|+|G|\le C\Lambda^{C_2}f.
 \label{eq:phys-residuals}
\end{equation}
Here and below the exponents of $\Lambda$ are fixed geometric
constants.

On the physical input plane the normalized variables satisfy
\begin{equation}
 (1+|A_*|^2)u=-T-y_*-2rA_*\cdot x-r^2w.
 \label{eq:phys-input-plane}
\end{equation}
The absolute Jacobian of $(x,w)\mapsto(z_\perp,z_3)$ is
\begin{equation}
 \frac{r^4}{1+|A_*|^2}.
 \label{eq:phys-slice-jacobian}
\end{equation}
Indeed $z_3=T+u$ and $z_\perp=b_*+uA_*+rx$; subtracting $A_*$
times the last row from the first two rows of their derivative
matrix leaves determinant $r^2\partial_wu$.
The crude volume of the bounded normalized box proves the lemma
for round tests.  It also proves it whenever $f$ is larger than a
sufficiently small fixed inverse power of $\Lambda$.

For the remaining $f$, the normalized coordinates are bounded by
constants, and Equation~\eqref{eq:phys-residuals} is as small as
needed.  The discriminant identity reads
\begin{equation}
 |F_x+2VF_w|^2=1+4(\beta+ku)G-4kF.
 \label{eq:phys-gradient-identity}
\end{equation}
It implies $|F_x+2VF_w|\ge1/2$.  Let $h(x,w)$ be the restriction
of $F$ to Equation~\eqref{eq:phys-input-plane}.  Put
\[
 a=\frac{2rA_*}{1+|A_*|^2},\qquad
 c_* =\frac{r^2}{1+|A_*|^2}.
\]
Then $h_x=F_x-aF_u$, $h_w=F_w-c_*F_u$, and
\begin{equation}
 G=\frac{1+|A_*+rV|^2}{1+|A_*|^2}F_u
               +V\cdot h_x+|V|^2h_w.
 \label{eq:phys-restricted-gradient}
\end{equation}
If $r$ is bounded by a fixed constant, a nonempty compatible set has bounded
$A_*$, because its actual row velocities $A=A_*+rV$ are bounded.
The coefficient of $F_u$ in
Equation~\eqref{eq:phys-restricted-gradient} is then bounded below.
Together with Equation~\eqref{eq:phys-gradient-identity}, this gives
$|\nabla h|\ge c>0$ once $|G|$ is sufficiently small.

For large $r$, replace the angular center by $0$.  This is a boost
of size $A_*/r$ in normalized coordinates; that size is bounded
because $|A_*|\le r+A_0+C_\psi$.  The boost preserves the discriminant
and changes all bounded coordinates and coefficients by bounded
amounts.  Now $|V|\le C/r$, while $|F_w|=|\beta+ku|\le C$.
Equation~\eqref{eq:phys-gradient-identity} gives $|F_x|\ge1/4$ for
all sufficiently large $r$.  In these coordinates $u$ in
Equation~\eqref{eq:phys-input-plane} is independent of $x$, so
$h_x=F_x$.  Choosing the cutoff between the two ranges once and for
all proves a uniform lower bound for the restricted gradient.

The polynomial $h$ has degree at most two.  In a fixed bounded
box, the set
$\{|h|\le q,\ |\nabla h|\ge c\}$ has volume at most $Cq$,
uniformly in its coefficients: partition according to a coordinate
derivative of magnitude at least $c/\sqrt3$, hold the other
coordinates fixed, and split the one-dimensional fiber at the
zeros of that derivative.  There are at most two monotonic pieces,
on each of which the mean value theorem gives length $O(q)$.
Applying this with $q=C\Lambda^{C_2}f$, and then using
Equation~\eqref{eq:phys-slice-jacobian} (or $r^4$ after recentering),
gives $|E|\le C\Lambda^{C_3}r^4f$ in the small-$f$ case.
Since $f\le f^{1-\eta}$, this suffices.  In the complementary
case the missing factor $f^{1-\eta}$ costs only another fixed power
of $\Lambda$, as already noted.
\end{proof}

\begin{proposition}[Regularity of the initial analyzed state]
\label{prop:phys-input-regularity}
Fix $0<\eta<1$ as in Definition~\ref{rt:regularity}. For every
$\gamma>0$, the state in
Equation~\eqref{eq:phys-round-state} has input regularity constant
\begin{equation}
 \kappa=C_{\eta,\gamma}H^{-\gamma}B^2
 \label{eq:phys-kappa}
\end{equation}
at each of the bounded starting times $T$.  Under
Equation~\eqref{eq:phys-energy-range}, it also satisfies the
polynomial-size and negligible-tail hypotheses of the round-wave
experiment with terminal length $s=H$.
\end{proposition}

\begin{proof}
Choose $\gamma_0>0$ so small that $C_1\gamma_0<\gamma/2$ in
Lemma~\ref{lem:phys-predecessor-volume}.  For a test $\mathcal T$,
restrict $g$ to its predecessor set $E$.  The truncated operator
of Lemma~\ref{lem:phys-initial-kernel} sends $g\mathbf1_{K\setminus E}$
to zero on $\mathcal T$.  The unrestricted operator on
$g\mathbf1_E$ has uniformly bounded $L^2$ norm by
Equation~\eqref{eq:phys-energy}.  Thus, for every $N$,
\[
 \mu_T(\mathcal T)
 \le C\norm{g\mathbf1_E}_2^2+C_NH^N\norm g_2^2
 \le C B^2H^{-C_1\gamma_0}r^4f^{1-\eta}+C_NH^N B^2.
\]
The uncertainty condition gives the uniform weight floor
\begin{equation}
 r^4f^{1-\eta}=(rf)^4f^{-3-\eta}\ge H^2.
 \label{eq:phys-weight-floor}
\end{equation}
Taking $N$ large proves regularity with
Equation~\eqref{eq:phys-kappa}, including every admissible width.

The ratios $e/\kappa$ and $\kappa/e$ are bounded above by fixed
powers of $H^{-1}$ by Equation~\eqref{eq:phys-energy-range}.
The kernel proof gives negligible phase tails outside a
fixed polynomial box: the input is compact, all velocities and
durations are bounded, and arbitrary decay is available.  Its
pointwise kernel bound also gives a phase-density bound by a fixed
power of $H^{-1}$ times $B^2$, and hence times $e$.  All these
estimates are uniform along any sequence of the present inputs.
\end{proof}

\subsection{Recovering the physical norm from leaf masses}

\begin{lemma}[Synthesis at the terminal scale]
\label{lem:phys-synthesis}
Let $v$ be the wave in Equation~\eqref{eq:phys-round-state}, and
let $\Omega$ be a fixed bounded subset of $(b,y)$ space.
At each left endpoint $T_j=T+jH$ in $[T,T+1]$, use a plateau
analysis of length $H$, and group its row masses into spatial
$H$-grid cubes.  Restrict to a fixed enlargement of $\Omega$, and
denote the resulting masses by $m_{j,Q}$.  They are admissible
round-test assignments with weights bounded above and below by
positive constants.  For every $0<\sigma<1$ and $N$,
\begin{equation}
 \int_T^{T+1}\int_\Omega |v|^3
 \le C_\sigma H^{-C_4\sigma}H^{-1/2}
                 \sum_{j,Q}m_{j,Q}^{3/2}
       +C_{N,\sigma}H^N e^{3/2}.
 \label{eq:phys-synthesis}
\end{equation}
Here $H$ is dyadic, so the time partition is exact.
\end{lemma}

\begin{proof}
At length $H$, the angular aperture is $d_H=1$ and both spatial
packet widths are $H$.  All row angles lie in a fixed bounded set.
A fixed round test with $r\ge1$, centered at a given cube, contains
that cube and all these angles.  Its weight $r^4$ is a fixed
constant, and different cube assignments are disjoint in row
space.

Synthesize the plateau analysis and propagate for $0\le h\le H$.
The same frequency calculation as in
Lemma~\ref{lem:phys-initial-kernel}, now with frequency boxes of
volume $O(H^{-3})$, gives a kernel bounded by
\[
 C_L H^{-3}\left(1+
       \frac{|(b,y)-(b',y')|}{H}\right)^{-L}.
\]
The bounded ray displacement during $h$ is $O(H)$ and has been
absorbed in this bound.  Its squared integral over source rows is
$O(H^{-3})$.  For an output point in a cube $Q$, retain only rows
in cubes within distance $C H^{1-\sigma}$ of $Q$.  Cauchy--Schwarz
then bounds the retained contribution by
\[
 C H^{-3/2}
 \left(\sum_{Q'\sim Q}m_{j,Q'}\right)^{1/2}.
\]
The omitted contribution is at most $C_{N,\sigma}H^N\sqrt e$
pointwise, by arbitrary kernel decay.  For $H$ sufficiently small,
all retained cubes lie in the chosen fixed enlargement of $\Omega$.
There are $O(H^{-3\sigma})$ neighbors of any $Q$, with the same
bound on overlap.  The spacetime cell over $Q$ has volume $H^4$.
Cubing the pointwise estimate, integrating over that cell, and
using
\[
 \left(\sum_{Q'\sim Q}m_{j,Q'}\right)^{3/2}
 \le \#\{Q'\sim Q\}^{1/2}
              \sum_{Q'\sim Q}m_{j,Q'}^{3/2}
\]
proves Equation~\eqref{eq:phys-synthesis}; for example, any fixed
$C_4\ge5$ suffices after harmless enlargement.  There are only
$O(H^{-4})$ spacetime cells, and arbitrary decay absorbs their
number in the error term.
\end{proof}

We explain explicitly how to use the exponent formulation of
Theorem~\ref{thm:round}.  Its wave bound, with regularity parameter
$\eta$, is $\Gamma_{\mathrm{wav},\mathrm{round}}\le100\eta$.  Hence for the
uniform class of experiments constructed above and every $\zeta>0$,
\begin{equation}
 H^{-1/2}\sum_{j,Q}m_{j,Q}^{3/2}
 \le C_{\eta,\gamma,\zeta}H^{-100\eta-\zeta}e\sqrt\kappa.
 \label{eq:phys-round-application}
\end{equation}
The bounded weights have been absorbed into the constant.
For completeness, at each fixed $H$ the normalized functional is
at most $C H^{-3/2}\sqrt{e/\kappa}\le C H^{-3/2}$: there are
$H^{-1}$ cuts and each has total assigned mass at most $e$.
Thus, if no constant in Equation~\eqref{eq:phys-round-application}
existed, there would be a sequence $H\downarrow0$ along
which the normalized functional exceeds
$H^{-100\eta-\zeta}$.  The regularity and tail estimates of
Proposition~\ref{prop:phys-input-regularity} give one
admissible polynomial sequence with all fixed setup constants
uniform.  Its upper growth exponent would be at least
$100\eta+\zeta$, contrary to Theorem~\ref{thm:round}.

Combining Equations~\eqref{eq:phys-synthesis},
\eqref{eq:phys-round-application}, and \eqref{eq:phys-kappa} gives
\begin{equation}
 \int_T^{T+1}\int_\Omega |v|^3
 \le C H^{-100\eta-\zeta-\gamma/2-C_4\sigma}
                    B\norm g_2^2.
 \label{eq:phys-level-estimate}
\end{equation}
The error term is absorbed here using
$e^{3/2}\le C B\norm g_2^2$.  The small-energy case was already
handled in Equation~\eqref{eq:phys-small-energy}.  Given any
$\beta>0$, first take $100\eta<\beta/4$, then
$\zeta<\beta/4$, $\gamma/2<\beta/4$, and
$C_4\sigma<\beta/4$; finally choose $\gamma_0$ as in
Proposition~\ref{prop:phys-input-regularity}.  Thus the loss in
Equation~\eqref{eq:phys-level-estimate} is smaller than any prescribed
$\beta$.  The finitely many null-time intervals and the constant
Jacobian in Equation~\eqref{eq:phys-null-coordinates} do not change
this conclusion.

\begin{proof}[Proof of Proposition~\ref{prop:phys-annular}]
It remains to pass from bounded amplitude pieces to arbitrary
localized input.  Normalize $\norm g_3=1$ with $\supp g\subset K$.
Fix an integer $M>4$.  The parts where $|g|<H^M$ and where
$|g|>H^{-M}$ have $L^1$ norms at most $|K|H^M$ and $H^{2M}$,
respectively.  The elementary bound
$\norm{K_{H,\tau}}_\infty\le C H^{-3}$ shows that their contributions
on the bounded output region are uniformly bounded, indeed tend
to zero.  No frequency restriction has been applied to $g$ before
making these spatial and amplitude restrictions.

Split the remaining part into disjoint dyadic amplitude pieces
$g_\nu$.  Their number is $L=O_M(1+\log H^{-1})$, and, with
$B_\nu=\norm{g_\nu}_\infty$,
\[
 B_\nu\norm{g_\nu}_2^2\le2\norm{g_\nu}_3^3,
 \qquad \sum_\nu\norm{g_\nu}_3^3\le1.
\]
Equation~\eqref{eq:phys-level-estimate}, with parameters making its
exponent less than $\delta$, and the inequality
\[
 \norm{\sum_{\nu=1}^L h_\nu}_3^3
 \le L^2\sum_\nu\norm{h_\nu}_3^3
\]
give the local cubed bound $C L^2H^{-\delta}$.
Since $L^2\le C_\delta H^{-\delta}$, taking a cube root yields a
loss at most $H^{-2\delta/3}\le H^{-\delta}$.  Homogeneity removes
the normalization.  Sum the cubed estimates over the output cubes,
using the bounded overlap of their enlarged input cubes and
Equation~\eqref{eq:phys-far-kernel}.  Finally sum the fixed angular
partition.  This proves Proposition~\ref{prop:phys-annular}.
\end{proof}

\subsection{Sobolev summation and the remaining exponents}

\begin{proof}[Proof of Theorem~\ref{thm:main}]
Let $h=J^\eps f$.  Use a smooth radial partition into a low-frequency
cutoff and annuli $|\xi|\asymp2^n$, $n\ge0$.  On the $n$th annulus
put the multiplier $J^{-\eps}$ into $a_H$, where $H=2^{-n}$.
After factoring out $H^\eps$, the rescaled symbols have uniformly
bounded derivatives.  Proposition~\ref{prop:phys-annular}, applied
with $\delta=\eps/2$, gives
\[
 \norm{W(P_nJ^{-\eps}h)}_{L^3_{x,\tau}(\R^3\times[1,2])}
 \le C_\eps 2^{-n\eps/2}\norm h_3.
\]
The series is summable by the triangle inequality.

We also record a direct justification at frequency zero.  For a
smooth cutoff $\chi$ supported in $|\xi|\le2$, decompose
\[
 \chi(\xi)e^{-i\tau|\xi|}
 =\chi(\xi)+\chi(\xi)(e^{-i\tau|\xi|}-1).
\]
The first term has an integrable convolution kernel.  Split the
second into annuli of radius $\lambda=2^{-j}\le1$.  In variables
$\xi=\lambda\omega$, its symbol and its derivatives through order
$4$ are $O(\lambda)$ on a fixed annulus, uniformly for
$1\le\tau\le2$.  Integration by parts using $(1-\Delta_\omega)^2$
shows that its inverse Fourier transform has $L^1$ norm
$O(\lambda)$.  These norms sum.  The smooth low-frequency factor
$J^{-\eps}$ preserves the same bounds.  Young's inequality
therefore handles the low-frequency term.  The resulting inequality
is exactly the negative-sign version of Theorem~\ref{thm:main};
complex conjugation proves the stated sign.
\end{proof}

\begin{corollary}[The positive-regularity local-smoothing range]
\label{cor:full-range}
For every $2<p<\infty$ and
$\alpha>\max\{0,1-3/p\}$,
\begin{equation}
 \norm{Uf}_{L^p(\R^3\times[1,2])}
 \le C_{p,\alpha}\norm{J^\alpha f}_{L^p(\R^3)}
 \qquad(f\in\mathcal S(\R^3)).
 \label{eq:phys-full-range}
\end{equation}
\end{corollary}

\begin{proof}
Take a smooth radial annular cutoff
$P_\lambda=\phi(|D|/\lambda)$, $\lambda\ge1$.  Plancherel gives
\[
 \norm{UP_\lambda}_{L^2_x\to L^2_{x,\tau}}\le C,
\]
and Proposition~\ref{prop:phys-annular} gives, for every $\delta>0$,
$\norm{UP_\lambda}_{L^3_x\to L^3_{x,\tau}}\le C_\delta\lambda^\delta$.
We first prove the remaining endpoint bound
\begin{equation}
 \norm{UP_\lambda}_{L^\infty_x\to L^\infty_{x,\tau}}
 \le C\lambda.
 \label{eq:phys-infinity-bound}
\end{equation}
The radial kernel, at $r=|x|>0$, is a fixed constant times
\[
 \frac1r\int_0^\infty e^{i\tau\rho}
             \rho\phi(\rho/\lambda)\sin(r\rho)\,d\rho.
\]
Put $I(s)=\int_0^\infty e^{is\rho}\rho\phi(\rho/\lambda)\,d\rho$.
For every $N$,
\[
 |I(s)|\le C_N\lambda^2(1+\lambda|s|)^{-N},\qquad
 |I'(s)|\le C_N\lambda^3(1+\lambda|s|)^{-N}.
\]
The kernel is a constant times
$r^{-1}(I(\tau+r)-I(\tau-r))$.  If $0<r\le1/2$, the mean value
theorem and $1\le\tau\le2$ bound it by $C_N\lambda^{3-N}$,
which also removes the apparent singularity at $r=0$.  If
$r\ge1/2$, its absolute value is at most
\[
 C_N\frac{\lambda^2}{r}
 \big((1+\lambda|\tau-r|)^{-N}
                +(1+\lambda(\tau+r))^{-N}\big).
\]
Multiplying by $r^2$ and integrating in $r$ gives $C\lambda$:
the first summand uses
$\int_{1/2}^\infty r(1+\lambda|\tau-r|)^{-N}\,dr\le C/\lambda$,
and the second is smaller.  Hence the kernel has $L^1$ norm at
most $C\lambda$, proving Equation~\eqref{eq:phys-infinity-bound}.

Interpolation between the $L^2$ and $L^3$ estimates gives an
arbitrarily small positive power of $\lambda$ for every $2<p\le3$.
For $p>3$, interpolate the $L^3$ bound with
Equation~\eqref{eq:phys-infinity-bound}; the $L^3$ interpolation
weight is $3/p$, giving
\[
 \norm{UP_\lambda}_{L^p_x\to L^p_{x,\tau}}
 \le C_{p,\delta}\lambda^{1-3/p+3\delta/p}.
\]
Choose $\delta$ small enough that this exponent, or its counterpart
for $p\le3$, is strictly below $\alpha$.  The same estimates hold
for uniformly smooth radial annular symbol families, as both the
kernel proof and Proposition~\ref{prop:phys-annular} are uniform
for those families.  We may thus insert $J^{-\alpha}$ into each
annular symbol, factor out $\lambda^{-\alpha}$, and sum the
resulting geometric series.  The low-frequency kernels have
uniformly bounded $L^1$ norms by the preceding proof, which applies
to every $L^p$.  This proves Equation~\eqref{eq:phys-full-range}.
\end{proof}

The same estimate describes the wave equation with independent initial
displacement $u_0$ and velocity $u_1$. For Schwartz data its solution is
\[
 u(t)=\cos(t|D|)u_0+\frac{\sin(t|D|)}{|D|}u_1,
\]
and for the exponents in Corollary~\ref{cor:full-range},
\[
 \norm{u}_{L^p(\R^3\times[1,2])}
 \le C_{p,\alpha}
 \bigl(\norm{J^\alpha u_0}_p+\norm{J^{\alpha-1}u_1}_p\bigr).
\]
At high frequency this follows from the two signs of the half-wave
estimate and the order $-1$ velocity multiplier. At low frequency,
retain the combined multiplier $\sin(t|\xi|)/|\xi|$: it is smooth
at zero and, after a smooth frequency cutoff, has a uniformly
integrable kernel for $1\le t\le2$. Density extends the solution
estimate to the indicated Sobolev data.

\subsection{Maximal half-waves and convergence to the initial data}

The spacetime estimate also controls the largest value of the wave over
a finite time interval, at the cost of additional regularity. This gives
a pointwise interpretation of the evolution for Sobolev initial data.
For $s\in\R$ and $1<p<\infty$, write
$W^{s,p}(\R^3)=\{f\in\mathcal S'(\R^3):J^s f\in L^p(\R^3)\}$,
with norm $\norm{J^s f}_p$.

\begin{corollary}[Maximal half-wave and pointwise convergence]
\label{cor:wave-maximal}
Let $3\le p<\infty$ and $s>1-2/p$. Then
\begin{equation}\label{eq:wave-maximal}
 \left\|\sup_{0<t<1}|Uf(\cdot,t)|\right\|_{L^p(\R^3)}
 \le C_{p,s}\norm{J^s f}_{L^p(\R^3)}
 \qquad(f\in\mathcal S(\R^3)).
\end{equation}
The family has a bounded maximal extension to $W^{s,p}(\R^3)$, and
in that extension every $f\in W^{s,p}(\R^3)$ satisfies
\begin{equation}\label{eq:wave-initial-limit}
 \lim_{t\downarrow0}Uf(x,t)=f(x)
 \quad\text{for almost every }x\in\R^3.
\end{equation}
\end{corollary}

\begin{proof}
For $p\ge3$, Corollary~\ref{cor:full-range} gives the $1/p-$
local smoothing gain relative to the fixed-time loss $1-2/p$.
The established maximal implication
\cite[Section~3.5, Equation~(3.24)]{BeltranHickmanSogge2019Survey}
therefore gives \eqref{eq:wave-maximal}. At each frequency, Sobolev
embedding in time costs arbitrarily more than $1/p$ derivatives;
together with local smoothing, this requires regularity arbitrarily
close to $1-2/p$ from above. The strict inequality on $s$ absorbs
the remaining positive losses. The implication uses scaling to cover
the time intervals approaching zero.

For Schwartz data, the Fourier integral converges uniformly to the
initial data as $t\downarrow0$. Schwartz functions are dense in
$W^{s,p}$, and the maximal bound gives the extension and carries this
convergence to every such datum. More explicitly, for a Schwartz
approximation $g$ to $f$, the measure of the set where
$\limsup_{t\downarrow0}|Uf-f|>\eta$ is at most
$C_{p,s}\eta^{-p}\norm{J^s(f-g)}_p^p$. Letting $g\to f$ proves
\eqref{eq:wave-initial-limit}.
\end{proof}
